\chapter{Schwinungen und Wellen}

\startitemize[packed]
\item Wellen treten in verschiedensten Formen auf: Wasserwellen, Schallwellen, elektromagnetische Wellen
\item Eine Welle ist eine räumliche und zeitliche Zustandsänderung physikalischer Größen, die meist nach bestimmten periodischen Gesetzmäßigkeiten erfolgt.
\item Die Ausbreitung einer Welle ist ein Energietransport, aber kein Materialtransport.
\stopitemize

\startdefinition[title={Welle}]
Eine Welle ist eine räumliche und zeitliche Zustandsänderung physikalischer Größen, die meist nach bestimmten periodischen Gesetzmäßigkeiten erfolgt[@LeifiWellen].
\stopdefinition

Die Ausbreitung {\bf mechanischer Wellen} erfordert einen Träger in dem sich
schwingungsfähige Teilchen befinden. Ein solcher Träger kann dabei fest,
flüssig oder auch gasförmig sein.

Weiter müssen diese schwingungsfähigen Teilchen untereinander eine Kopplung
aufweisen, so dass sich die von außen einwirkende Störung über das System
fortpflanzen kann.

Ein Erreger zwingt ein Teilchen aus seiner {\em Ruhelage}\sidenote{{\em Ruhelage} ist die \quotation{Gleichgewichtslage}}. Aufgrund der Trägheit
übernimmt das nächte Teilchen die Störung etwas zeitversetzt und so entsteht
eine Phasenverschiebung $\Delta \phi$ zwischen den Bewegungen benachbarter
Teilchen. Auf diese Weise pflanzt sich die Störung durch den Körper fort.

\section{Überlagerung von Wellen}

Wellen können miteinander Interferieren. Dabei unterscheidet man zwischen
\quotation{konstruktiver} (Versterkender) und
\quotation{destruktiver}\sidenote{destruktive Interferenz findet auch häufige
anwendung im Täglichen Lebene. Zum Beispiel funktioniert \quotation{active
noice cancelling} auf diese art und weise.} (Schwächender) Interferenz. Bei
konstruktiver Interferenz treffen ein Tal auf ein Tal, oder ein Berg auf einen
Berg (\in{s. Abb.}[fig:waves_no_offset]). Bei destruktiver Interferenz Treffen
ein Tal auf einen Berg (\in{s. Abb.}[fig:waves_offset]).

\placefigure[here][fig:waves_no_offset]{Konstruktiver Interferenz zweier identischer Wellen. Hier zwei mal $\sin(2x)$, dargestellt in lila.}{
\startMPcode
	u := 1cm;

    draw function (1, "x", ".5 * sin(2x)", 1, 10, .01) scaled u
    withpen pencircle scaled 1mm withcolor transparent(1,.5,red);

    draw function (2, "x", ".5 * sin(2x)", 1, 10, .01) scaled u
    withpen pencircle scaled 1mm withcolor transparent(1,.5,blue);

    draw function (3, "x", ".5 * (sin(2x)+sin(2x))", 1, 10, .01) scaled u
    withpen pencircle scaled 1mm withcolor transparent(1,.5,black);
\stopMPcode
}

\placefigure[here][fig:waves_offset]{Destruktive Interferenz zweier versetzter Welen, hier $\sin(2x + \pi)$ (rot) und $\sin(2x)$ (blau)}{
\startMPcode
	u := 1cm;

    draw function (1, "x", "sin(2x + pi)", 1, 10, .01) scaled u
    withpen pencircle scaled 1mm withcolor transparent(1,.5,red);

    draw function (2, "x", "sin(2x)", 1, 10, .01) scaled u
    withpen pencircle scaled 1mm withcolor transparent(1,.5,blue);

    draw function (3, "x", "sin(2x)+sin(2x + pi)", 1, 10, .01) scaled u
    withpen pencircle scaled 1mm withcolor transparent(1,.5,black);
\stopMPcode
}

\section{Lidar vs Radar}


Lidar, Radar, oder auch Sonar, funktionieren alle sehr ähnlich, indem
Energiewellen ausgesendet werden und auf eine Reflexion gewartet wird.

\subsection[title={Radar},reference={radar}]

\quotation{{\bf Ra}dio {\bf d}etection {\bf a}nd {\bf r}anging}, kurz
Radar, ist die Bezeichnung für verschiedene Erkennungs- oder auch
Ortungsverfahren, welche mit {\bf elektromagnetischen Wellen} im
Radiofrequenzbereich (kleiner als
$300\text{MHz}$\footnote{\cite[alternative=entry][Radiowellen_LEIFI]})

\subsubsection[title={Funktionsweise},reference={funktionsweise}]

Für die Funktionsweise von Radar-Sensoren sind drei grundlegende
physikalische Gesetze von
bedeutung\footnote{\cite[alternative=entry][Wolff2024Aug]}:

\startitemize[packed]
\item
  Die {\bf Reflexion} elektromagnetischer Wellen.
\item
  Die {\bf konstante Ausbreitungsgeschwindigkeit} der
  elektromagnetischen Wellen, welche sich mit knapp Lichtgeschwindigkeit
  ausbreiten (ca $300 \frac{km}{s}$).
\item
  Die {\bf geradlinige Ausbreitung} der elektromagnetischen Wellen.
\stopitemize

Die elektromagnetischen Wellen werden also ausgesendet und deren
Reflexion gemessen. Mit der Zeitdifferenz des aussenden und des
empfangen kann die Distanz ermittelt werden, welche die Wellen
zurückgelegt haben.

\subsubsection[title={Anwendung},reference={anwendung}]

Die Anwendungsmöglichkeiten sind vielfältig. Die Technologie wird unter
anderem dafür genutzt, um Positionen zu bestimmen, oder auch füllstände
in Lagertanks\footnote{\cite[alternative=entry][IFM_Radar]}

\subsection[title={Lidar},reference={lidar}]

„{\bf Li}ght {\bf d}etection {\bf a}nd {\bf r}anging“, kurz Lidar, ist
eine Technologie, welche die präziese messung von Entfernungen mit sehr
geringer latenz erlaubt\footnote{\cite[alternative=entry][IBM_LIDAR]}.

\subsubsection[title={Funktionsweise},reference={funktionsweise-1}]

Klassische Lidar-Sensoren bezitzen einen Transmitter und einen
Empfänger. Über den Transmitter wird ein {\bf gepulster Laser} wird
ausgesendet. Der Zeitpunkt dessen Reflexion wird mit dem Empfänger
gemessen und somit die \quotation{Flugdauer} ermittelt. Da die
Lichtgeschwindigkeit konstant ist (ca $300
\frac{km}{s}$\footnote{\cite[alternative=entry][lichtgeschwindigkeit_LEIFI]})
ist diese Methode sehr schnell und präziese.


\subsubsection[title={Anwendung},reference={anwendung-1}]

Lidar Technologie findet in der heutigen Welt eine breite und
vielschichtige Anwendung. Von Forschung bis Technologie im Alltag,
Lidar-Sensoren sind überall zu finden und auch nicht mehr weg zu denken.

In der Archäologie wurde 2018 mithilfe von Lidar-Luftscans eine bisher
übersehene Maya-Ruinenstadt gefunden. Dank der Scans konnte eine
großflächige und detailreiche 3D-Karte kreiert werden, welche die
Umrisse der Ruinen
aufdeckte\footnote{\cite[alternative=entry][lidar_defunk]}.

Auch für das autonome Fahren spielt die Lidar-Technologie eine große
Rolle. Die Autos müssen ihre Umgebung in drei Dimensionen erfassen
können. Die Lidar-Sensoren werden mit vielen anderen Sensoren, wie auch
Kameras, kombiniert um die aktuelle Situation einzuschätzen. Lidar ist
deshalb geeignet, da es sehr präzise Messungen mit sehr geringer Latenz
ermöglicht\footnote{\cite[alternative=entry][9127855]}. Aber nicht nur
für das Autonome steuern von Autos wird diese Technologie genutzt, auch
Staubsaug-Roboter nutzen sie, um sich in einem Raum zurecht zu finden.


\subsection{Funktionsweise}

\startframedtask
Erklären Sie die Funktionsweise dieser beiden Technologien physikalisch.
Geben Sie dafür auch die Gleichungen an, mit denen die nötigen
Berechnungen durchgeführt werden können. Verzichten Sie auf Erklärungen
zur Nutzung der Raman-Streuung mit LIDAR-Systemen

\stopframedtask

Die Messung der \quotation{Time of flight}, also der Strecke, die die
Energiewelle in welcher Zeit zurückgelegt hat, sieht wie folgt
aus\footnote{\cite[alternative=entry][Gotzig2015Mar]}:

\startformula 
s=\frac{c_0\cdot t}{2}
 \stopformula

\startitemize[packed]
\item
  $s =$ Abstand
\item
  $c_0 =$ Geschwindigkeit der Energiewelle, bei Lidar:
  Lichtgeschwindigkeit (ca. $300 \frac{km}{s}$) und bei Radar:
  geschwindigkeit ausbreitung von elektromagnetischer Wellen (annähernd
  Lichtgeschwindigkeit)
\item
  $t =$ Zeit in $s$.
\stopitemize

\startformula 
\lambda = c \cdot T = \frac{c}{f}
 \stopformula


\subsection{Einordnung}

\startplacetable[location=none]
\setupTABLE[r][1][style=bold,background=color,backgroundcolor=gray,width=2.5cm]
\setupTABLE[c][1][style=bold,background=color,backgroundcolor=gray]
\bTABLE
\startCSV
[background=DiagonalRule] \DiagonalLabel{Kategorie}{Art}, Radar, Lidar
Wellenlänge:, $1m$, zwischen $1m$ und $780nm$
Frequenzen:, $300MHz$, $300GHz$ bis $385THz$
Art:, Radiowellen, Infrarot
\stopCSV
\eTABLE
\stopplacetable
